% GENERATED FILE. Edit canonical lessons, then run npm run generate:books.
% canonical-source-hash: fnv1a64:f11e0c3ac70e400c
% canonical-lessons: MR-C139-tekdi, MR-C139-jhara, MR-C139-dhabdhaba, MR-C139-gavat, MR-C139-phandi

\chapter{Hill, Spring, Waterfall, Grass, Branch}
\label{ch:mr-hill-spring-waterfall-grass-branch}
\hlchaptermodality{139}

\begin{chapteropening}
\textbf{By the end of this chapter:} \emph{I can say a hill, a spring, a waterfall, grass and a branch.}

Complete the last lesson of chapter 139: i can say a hill, a spring, a waterfall, grass and a branch.
\end{chapteropening}

\section[ṭekḍī]{\mr{टेकडी} (ṭekḍī) — a hill}

\label{lesson:MR-C139-tekdi}



\begin{quote}
Before the new one: say the Marathi for a coconut, then the Marathi for flour.
\end{quote}

\subsection*{You'll want to know: \mr{टेकडी}}
\textbf{\mr{टेकडी}} — \emph{ṭekḍī} — \textquotedblleft{}a hill\textquotedblright{}.

\begin{grammarlens}[title={What you've built}]
One more word for this chapter.
\end{grammarlens}

\subsection*{Your turn}
\begin{itemize}
\raggedright
  \item \emph{Say it:} \emph{ṭekḍī}
  \item \emph{Say it:} \emph{ṭekḍī}, once more
  \item \emph{Say it:} say \emph{ṭekḍī}
  \item \emph{Recall:} say the Marathi for a coconut, then the Marathi for flour, then say \emph{ṭekḍī} again
\end{itemize}

\begin{tcolorbox}[breakable,colback=teal!4,colframe=teal!35!black,title={Before you move on}]
What does \mr{टेकडी} mean? (\textquotedblleft{}A hill\textquotedblright{}.) Say it once more.
\end{tcolorbox}
\section[jharā]{\mr{झरा} (jharā) — a spring, a stream}

\label{lesson:MR-C139-jhara}



\begin{quote}
Before the new one: say the Marathi for flour, then the Marathi for a hill.
\end{quote}

\subsection*{You'll want to know: \mr{झरा}}
\textbf{\mr{झरा}} — \emph{jharā} — \textquotedblleft{}a spring, a stream\textquotedblright{}.

\begin{grammarlens}[title={What you've built}]
One more word for this chapter.
\end{grammarlens}

\subsection*{Your turn}
\begin{itemize}
\raggedright
  \item \emph{Say it:} \emph{jharā}
  \item \emph{Say it:} \emph{jharā}, once more
  \item \emph{Say it:} say \emph{jharā}
  \item \emph{Recall:} say the Marathi for flour, then the Marathi for a hill, then say \emph{jharā} again
\end{itemize}

\begin{tcolorbox}[breakable,colback=teal!4,colframe=teal!35!black,title={Before you move on}]
What does \mr{झरा} mean? (\textquotedblleft{}A spring, a stream\textquotedblright{}.) Say it once more.
\end{tcolorbox}
\section[dhabdhabā]{\mr{धबधबा} (dhabdhabā) — a waterfall}

\label{lesson:MR-C139-dhabdhaba}



\begin{quote}
Before the new one: say the Marathi for a hill, then the Marathi for a spring.
\end{quote}

\subsection*{You'll want to know: \mr{धबधबा}}
\textbf{\mr{धबधबा}} — \emph{dhabdhabā} — \textquotedblleft{}a waterfall\textquotedblright{}.

\begin{grammarlens}[title={What you've built}]
One more word for this chapter.
\end{grammarlens}

\subsection*{Your turn}
\begin{itemize}
\raggedright
  \item \emph{Say it:} \emph{dhabdhabā}
  \item \emph{Say it:} \emph{dhabdhabā}, once more
  \item \emph{Say it:} say \emph{dhabdhabā}
  \item \emph{Recall:} say the Marathi for a hill, then the Marathi for a spring, then say \emph{dhabdhabā} again
\end{itemize}

\begin{tcolorbox}[breakable,colback=teal!4,colframe=teal!35!black,title={Before you move on}]
What does \mr{धबधबा} mean? (\textquotedblleft{}A waterfall\textquotedblright{}.) Say it once more.
\end{tcolorbox}
\section[gavat]{\mr{गवत} (gavat) — grass}

\label{lesson:MR-C139-gavat}



\begin{quote}
Before the new one: say the Marathi for a spring, then the Marathi for a waterfall.
\end{quote}

\subsection*{You'll want to know: \mr{गवत}}
\textbf{\mr{गवत}} — \emph{gavat} — \textquotedblleft{}grass\textquotedblright{}.

\begin{grammarlens}[title={What you've built}]
One more word for this chapter.
\end{grammarlens}

\subsection*{Your turn}
\begin{itemize}
\raggedright
  \item \emph{Say it:} \emph{gavat}
  \item \emph{Say it:} \emph{gavat}, once more
  \item \emph{Say it:} say \emph{gavat}
  \item \emph{Recall:} say the Marathi for a spring, then the Marathi for a waterfall, then say \emph{gavat} again
\end{itemize}

\begin{tcolorbox}[breakable,colback=teal!4,colframe=teal!35!black,title={Before you move on}]
What does \mr{गवत} mean? (\textquotedblleft{}Grass\textquotedblright{}.) Say it once more.
\end{tcolorbox}
\section[phāndī]{\mr{फांदी} (phāndī) — a branch}

\label{lesson:MR-C139-phandi}



\begin{quote}
Before the new one: say the Marathi for a waterfall, then the Marathi for grass.
\end{quote}

\subsection*{You'll want to know: \mr{फांदी}}
\textbf{\mr{फांदी}} — \emph{phāndī} — \textquotedblleft{}a branch\textquotedblright{}.

\begin{grammarlens}[title={What you've built}]
One more word for this chapter.
\end{grammarlens}

\subsection*{Your turn}
\begin{itemize}
\raggedright
  \item \emph{Say it:} \emph{phāndī}
  \item \emph{Say it:} \emph{phāndī}, once more
  \item \emph{Say it:} say \emph{phāndī}
  \item \emph{Recall:} say the Marathi for a waterfall, then the Marathi for grass, then say \emph{phāndī} again
\end{itemize}

\begin{tcolorbox}[breakable,colback=teal!4,colframe=teal!35!black,title={Before you move on}]
What does \mr{फांदी} mean? (\textquotedblleft{}A branch\textquotedblright{}.) Say it once more.
\end{tcolorbox}
